% Part: first-order-logic
% Chapter: completeness
% Section: completeness-theorem

\documentclass[../../../include/open-logic-section]{subfiles}

\begin{document}

\iftag{FOL}
      {\olfileid{fol}{com}{cth}}
      {\olfileid{pl}{com}{cth}}

\olsection{De volledigheidsstelling}

\begin{explain}
Door onze resultaten te combineren verkrijgen we de volledigheidsstelling.
\end{explain}

\begin{thm}[Volledigheidsstelling]
\ollabel{thm:completeness}
Laat $\Gamma$ een verzameling !!{sentence}s zijn. Als $\Gamma$ consistent
is, dan is zij vervulbaar.
\end{thm}

\begin{proof}
Veronderstel dat $\Gamma$ consistent is.\iftag{FOL}{ Volgens
  \olref[hen]{lem:henkin} bestaat er een verzadigde consistente
  verzameling $\Gamma' \supseteq \Gamma$.}{} Volgens
\olref[lin]{lem:lindenbaum} bestaat er een
$\Gamma^* \supseteq \iftag{FOL}{\Gamma'}{\Gamma}$ die consistent en
!!{complete} is.
\iftag{FOL}{
  Omdat $\Gamma' \subseteq \Gamma^*$, geldt voor elke
  !!{formula}~$!A(x)$ dat $\Gamma^*$ !!a{sentence} van de volgende vorm bevat:
  \iftag{prvEx}
        {$\lexists[x][!A(x)] \lif !A(c)$}
        {$\lnot\lforall[x][!A(x)] \lif \lnot !A(c)$}
  en dus is $\Gamma^*$ verzadigd. Als $\Gamma$ geen~$\eq$ bevat, volgt
  uit}{Uit}
\olref[mod]{lem:truth},
$\iftag{FOL}{\Sat{M(\Gamma^*)}}{\pSat{v(\Gamma^*)}}{!A}$ dan en slechts
dan als $!A \in \Gamma^*$. Hieruit volgt in het bijzonder dat voor alle
$!A \in \Gamma$, $\iftag{FOL}{\Sat{M(\Gamma^*)}}{\pSat{v(\Gamma^*)}}{!A}$,
zodat $\Gamma$ vervulbaar is.\iftag{FOL}{ Als $\Gamma$ wel~$\eq$ bevat,
  dan geldt volgens \olref[ide]{lem:truth} voor alle !!{sentence}s~$!A$
  dat $\Sat{\equivclass{M}{\approx}}{!A}$ dan en slechts dan als
  $!A \in \Gamma^*$. In het bijzonder geldt
  $\Sat{\equivclass{M}{\approx}}{!A}$ voor alle $!A \in \Gamma$, zodat
  $\Gamma$ vervulbaar is.}{}
\end{proof}

\begin{cor}[Volledigheidsstelling, tweede formulering]
\ollabel{cor:completeness}
Voor alle $\Gamma$ en !!{sentence}s~$!A$ geldt: als
$\Gamma \Entails !A$, dan $\Gamma \Proves !A$.
\end{cor}

\begin{proof}
Merk op dat de $\Gamma$'s in \olref{cor:completeness} en
\olref{thm:completeness} universeel gekwantificeerd zijn. Om verwarring
te voorkomen, formuleren we \olref{thm:completeness} opnieuw met een
andere variabele: voor elke verzameling !!{sentence}s~$\Delta$ geldt dat
als $\Delta$ consistent is, zij vervulbaar is. Door contrapositie geldt
dus: als $\Delta$ niet vervulbaar is, dan is $\Delta$ inconsistent. Dit
gebruiken we om het corollarium te bewijzen.

Veronderstel dat $\Gamma \Entails !A$. Volgens
\olref[syn][sem]{prop:entails-unsat} is $\Gamma \cup \{\lnot !A\}$ dan
onvervulbaar. Als we $\Gamma \cup \{\lnot !A\}$ als onze $\Delta$ nemen,
volgt uit de vorige formulering van \olref{thm:completeness} dat
$\Gamma \cup \{\lnot !A\}$ inconsistent is. Volgens
\iftag{FOL}{%
  \tagrefs{prfAX/{fol:axd:prv:prop:prov-incons},
    prfSC/{fol:seq:prv:prop:prov-incons},
    prfND/{fol:ntd:prv:prop:prov-incons},
    prfTab/{fol:tab:prv:prop:prov-incons}}}{%
  \tagrefs{prfAX/{pl:axd:prv:prop:prov-incons},
    prfSC/{pl:seq:prv:prop:prov-incons},
    prfND/{pl:ntd:prv:prop:prov-incons},
    prfTab/{pl:tab:prv:prop:prov-incons}}},
$\Gamma \Proves !A$.
\end{proof}

\tagprob{FOL}
\begin{prob}
Gebruik \olref[fol][com][cth]{cor:completeness} om
\olref[fol][com][cth]{thm:completeness} te bewijzen en toon daarmee aan
dat de twee formuleringen van de volledigheidsstelling equivalent zijn.
\end{prob}
\tagendprob

\tagprob{notFOL}
\begin{prob}
Gebruik \olref[pl][com][cth]{cor:completeness} om
\olref[pl][com][cth]{thm:completeness} te bewijzen en toon daarmee aan
dat de twee formuleringen van de volledigheidsstelling equivalent zijn.
\end{prob}
\tagendprob

\tagprob{FOL}
\begin{prob}
Om een !!{derivation}ssysteem volledig te laten zijn, moeten zijn regels
sterk genoeg zijn om elke onvervulbare verzameling inconsistent te
bewijzen. Welke regels van !!{derivation} waren nodig om de volledigheid
te bewijzen? Zijn sommige van deze regels nergens in het bewijs gebruikt?
Maak ter beantwoording een lijst of diagram waaruit blijkt welke regels
van !!{derivation} zijn gebruikt in welke resultaten die tot het bewijs
van \olref[fol][com][cth]{thm:completeness} leiden. Vermeld ook elk
stilzwijgend gebruik van regels in deze bewijzen.
\end{prob}
\tagendprob

\tagprob{notFOL}
\begin{prob}
Om een !!{derivation}ssysteem volledig te laten zijn, moeten zijn regels
sterk genoeg zijn om elke onvervulbare verzameling inconsistent te
bewijzen. Welke regels van !!{derivation} waren nodig om de volledigheid
te bewijzen? Zijn sommige van deze regels nergens in het bewijs gebruikt?
Maak ter beantwoording een lijst of diagram waaruit blijkt welke regels
van !!{derivation} zijn gebruikt in welke resultaten die tot het bewijs
van \olref[pl][com][cth]{thm:completeness} leiden. Vermeld ook elk
stilzwijgend gebruik van regels in deze bewijzen.
\end{prob}
\tagendprob

\end{document}
